% Adaptation de la slide 55 (page PDF 64) de la soutenance du 8 septembre 2026.
% Le constat vient de nos essais ; la redondance est une explication envisagée.
\begin{frame}[t]{Hessienne, gradient : le guidage n'a pas aidé}
  \begin{center}
    \resizebox{0.48\textwidth}{!}{\input{figs/lora_adaptation}}\par\smallskip
    {\small \textbf{Base : SAM gelé + LoRA} \citb{hu22}, comme dans CrackSAM \citb{ge24}.}
  \end{center}
  \vspace{0.10cm}
  \small
  \begin{itemize}
    \setlength{\itemsep}{7pt}
    \item \textbf{Nos essais :} ajouter la hessienne, le gradient ou Frangi
      \textbf{n'améliore pas} SAM + LoRA.
    \item \textbf{Explication possible :} ces informations locales semblent
      \textbf{déjà représentées par SAM}.
    \item \textbf{Limite de Frangi :} les ombres et la granularité
      provoquent des \textbf{sur-détections}.
  \end{itemize}
  \vspace{0.16cm}
  \begin{center}
    \formulebox{\parbox{0.83\textwidth}{\centering\small
      \textbf{Piste :} relier les zones d'une même fissure par une \textbf{hiérarchie}.}}
  \end{center}
  \reffoot{hu22,ge24}
\end{frame}
